\documentclass[11pt,a4paper]{article}

\usepackage[utf8]{inputenc}
\usepackage[T2A]{fontenc}
\usepackage[russian,english]{babel}
\usepackage{geometry}
\usepackage{enumitem}
\usepackage{booktabs}
\usepackage{array}
\usepackage{longtable}
\usepackage{xcolor}
\usepackage{titlesec}
\usepackage{hyperref}
\usepackage{fancyhdr}
\usepackage{tcolorbox}
\tcbuselibrary{breakable,listings}

\geometry{left=25mm, right=25mm, top=30mm, bottom=25mm}

% Colors
\definecolor{coilblue}{HTML}{1A5276}
\definecolor{coilgray}{HTML}{555555}
\definecolor{codebg}{HTML}{F4F4F4}
\definecolor{codeborder}{HTML}{CCCCCC}
\definecolor{errorbg}{HTML}{FFF0F0}

% Section styles
\titleformat{\section}{\Large\bfseries\color{coilblue}}{}{0em}{}[\vspace{-0.5em}\textcolor{codeborder}{\rule{\textwidth}{0.4pt}}]
\titleformat{\subsection}{\large\bfseries\color{coilblue}}{}{0em}{}
\titleformat{\subsubsection}{\normalsize\bfseries}{}{0em}{}

% Code environment (tcolorbox + listings for proper background fill)
\newtcblisting{code}{
  colback=codebg,
  colframe=codeborder,
  boxrule=0.4pt,
  arc=0pt,
  left=8pt, right=8pt, top=6pt, bottom=6pt,
  breakable,
  listing only,
  listing options={
    basicstyle=\small\ttfamily,
    breaklines=true,
    columns=fullflexible,
    keepspaces=true,
    literate={а}{{\cyra}}1 {б}{{\cyrb}}1 {в}{{\cyrv}}1 {г}{{\cyrg}}1
             {д}{{\cyrd}}1 {е}{{\cyre}}1 {ё}{{\cyryo}}1 {ж}{{\cyrzh}}1
             {з}{{\cyrz}}1 {и}{{\cyri}}1 {й}{{\cyrishrt}}1 {к}{{\cyrk}}1
             {л}{{\cyrl}}1 {м}{{\cyrm}}1 {н}{{\cyrn}}1 {о}{{\cyro}}1
             {п}{{\cyrp}}1 {р}{{\cyrr}}1 {с}{{\cyrs}}1 {т}{{\cyrt}}1
             {у}{{\cyru}}1 {ф}{{\cyrf}}1 {х}{{\cyrh}}1 {ц}{{\cyrc}}1
             {ч}{{\cyrch}}1 {ш}{{\cyrsh}}1 {щ}{{\cyrshch}}1 {ъ}{{\cyrhrdsn}}1
             {ы}{{\cyrery}}1 {ь}{{\cyrsftsn}}1 {э}{{\cyrerev}}1 {ю}{{\cyryu}}1
             {я}{{\cyrya}}1
             {А}{{\CYRA}}1 {Б}{{\CYRB}}1 {В}{{\CYRV}}1 {Г}{{\CYRG}}1
             {Д}{{\CYRD}}1 {Е}{{\CYRE}}1 {Ё}{{\CYRYO}}1 {Ж}{{\CYRZH}}1
             {З}{{\CYRZ}}1 {И}{{\CYRI}}1 {Й}{{\CYRISHRT}}1 {К}{{\CYRK}}1
             {Л}{{\CYRL}}1 {М}{{\CYRM}}1 {Н}{{\CYRN}}1 {О}{{\CYRO}}1
             {П}{{\CYRP}}1 {Р}{{\CYRR}}1 {С}{{\CYRS}}1 {Т}{{\CYRT}}1
             {У}{{\CYRU}}1 {Ф}{{\CYRF}}1 {Х}{{\CYRH}}1 {Ц}{{\CYRC}}1
             {Ч}{{\CYRCH}}1 {Ш}{{\CYRSH}}1 {Щ}{{\CYRSHCH}}1 {Ъ}{{\CYRHRDSN}}1
             {Ы}{{\CYRERY}}1 {Ь}{{\CYRSFTSN}}1 {Э}{{\CYREREV}}1 {Ю}{{\CYRYU}}1
             {Я}{{\CYRYA}}1
  }
}

% Error code environment (for anti-patterns)
\newtcblisting{badcode}{
  colback=errorbg,
  colframe=red,
  boxrule=0.4pt,
  arc=0pt,
  left=8pt, right=8pt, top=6pt, bottom=6pt,
  breakable,
  listing only,
  listing options={
    basicstyle=\small\ttfamily,
    breaklines=true,
    columns=fullflexible,
    keepspaces=true,
    literate={а}{{\cyra}}1 {б}{{\cyrb}}1 {в}{{\cyrv}}1 {г}{{\cyrg}}1
             {д}{{\cyrd}}1 {е}{{\cyre}}1 {ё}{{\cyryo}}1 {ж}{{\cyrzh}}1
             {з}{{\cyrz}}1 {и}{{\cyri}}1 {й}{{\cyrishrt}}1 {к}{{\cyrk}}1
             {л}{{\cyrl}}1 {м}{{\cyrm}}1 {н}{{\cyrn}}1 {о}{{\cyro}}1
             {п}{{\cyrp}}1 {р}{{\cyrr}}1 {с}{{\cyrs}}1 {т}{{\cyrt}}1
             {у}{{\cyru}}1 {ф}{{\cyrf}}1 {х}{{\cyrh}}1 {ц}{{\cyrc}}1
             {ч}{{\cyrch}}1 {ш}{{\cyrsh}}1 {щ}{{\cyrshch}}1 {ъ}{{\cyrhrdsn}}1
             {ы}{{\cyrery}}1 {ь}{{\cyrsftsn}}1 {э}{{\cyrerev}}1 {ю}{{\cyryu}}1
             {я}{{\cyrya}}1
             {А}{{\CYRA}}1 {Б}{{\CYRB}}1 {В}{{\CYRV}}1 {Г}{{\CYRG}}1
             {Д}{{\CYRD}}1 {Е}{{\CYRE}}1 {Ё}{{\CYRYO}}1 {Ж}{{\CYRZH}}1
             {З}{{\CYRZ}}1 {И}{{\CYRI}}1 {Й}{{\CYRISHRT}}1 {К}{{\CYRK}}1
             {Л}{{\CYRL}}1 {М}{{\CYRM}}1 {Н}{{\CYRN}}1 {О}{{\CYRO}}1
             {П}{{\CYRP}}1 {Р}{{\CYRR}}1 {С}{{\CYRS}}1 {Т}{{\CYRT}}1
             {У}{{\CYRU}}1 {Ф}{{\CYRF}}1 {Х}{{\CYRH}}1 {Ц}{{\CYRC}}1
             {Ч}{{\CYRCH}}1 {Ш}{{\CYRSH}}1 {Щ}{{\CYRSHCH}}1 {Ъ}{{\CYRHRDSN}}1
             {Ы}{{\CYRERY}}1 {Ь}{{\CYRSFTSN}}1 {Э}{{\CYREREV}}1 {Ю}{{\CYRYU}}1
             {Я}{{\CYRYA}}1
  }
}

% Inline code
\newcommand{\ic}[1]{\texttt{\small #1}}

% Header/footer
\pagestyle{fancy}
\fancyhf{}
\fancyhead[L]{\small\color{coilgray}COIL --- Cognitive Orchestration Interface Language}
\fancyhead[R]{\small\color{coilgray}\thepage}
\fancyfoot[L]{\small\color{coilgray}Animata Systems}
\fancyfoot[R]{\small\color{coilgray}Draft 0.2}
\renewcommand{\headrulewidth}{0.4pt}
\renewcommand{\footrulewidth}{0.4pt}

\begin{document}

% ============================================================
% TITLE
% ============================================================

\begin{center}
{\Huge\bfseries\color{coilblue} COIL}\\[6pt]
{\Large Cognitive Orchestration Interface Language}\\[12pt]
{\large\color{coilgray} Операторский язык сценариев для когнитивной оркестрации}\\[6pt]
{\color{coilgray} Учебник}\\[6pt]
{\small\color{coilgray} Черновик 0.2 --- март 2026}
\end{center}

\vspace{1em}

Этот учебник проведёт вас от первого протокола до трёх самых востребованных сценариев. По пути вы узнаете, какие решения принимает язык за вас, а какие остаются за вами~--- и как избежать типичных ошибок.

\tableofcontents
\newpage

% ============================================================
% SECTION 1: С ЧЕГО НАЧИНАТЬ
% ============================================================

\section{С чего начинать}

\subsection{Что такое протокол COIL}

Протокол COIL~--- это сценарий для агента: что обдумать, кого спросить, когда подождать, когда ответить. Он описывает один конкретный случай работы, а не «всё поведение» агента навсегда.

Протокол читается сверху вниз, как чек-лист. Каждая строка~--- явное действие: подумать, отправить сообщение, дождаться ответа, завершить. Никакой скрытой логики внутри длинного промпта.

\subsection{Минимальный протокол}

Самый короткий рабочий протокол состоит из четырёх шагов: подумать, подождать, написать, выйти.

\begin{code}
ДУМАЙ приветствие
  ЦЕЛЬ <<
  Придумай короткое приветственное сообщение.
  >>
  РЕЗУЛЬТАТ
  * текст: ТЕКСТ - приветствие
КОНЕЦ

ЖДИ
  НА ?приветствие
КОНЕЦ

НАПИШИ
  КОМУ @пользователь
  <<
  $приветствие.текст
  >>
КОНЕЦ

ВЫХОД
\end{code}

Разберём по шагам:

\begin{enumerate}[leftmargin=2em]
\item \ic{ДУМАЙ приветствие}~--- запускает когнитивный шаг. LLM получает задание и начинает думать. Оператор не ждёт результата~--- он сразу создаёт \textbf{обещание} \ic{?приветствие} и продолжает.

\item \ic{ЖДИ НА ?приветствие}~--- точка синхронизации. Протокол останавливается, пока LLM не закончит. После этого результат доступен как \textbf{значение} \ic{\$приветствие}.

\item \ic{НАПИШИ КОМУ @пользователь}~--- отправляет сообщение. Внутри шаблона \ic{<{<} >{>}} подставляется \ic{\$приветствие.текст}~--- поле из структурированного результата.

\item \ic{ВЫХОД}~--- завершает протокол.
\end{enumerate}

\subsubsection{Два главных сигила}

\begin{itemize}[leftmargin=2em]
\item \ic{\$имя}~--- \textbf{значение}. Данные, которые уже доступны.
\item \ic{?имя}~--- \textbf{обещание}. Результат, который ещё не готов.
\end{itemize}

Правило простое: запускающий оператор (\ic{ДУМАЙ}, \ic{ВЫПОЛНИ}, \ic{НАПИШИ}) создаёт \ic{?имя}. После \ic{ЖДИ} обещание разрешается в \ic{\$имя}.

\subsection{Добавляем контекст: роль и входные данные}

Минимальный протокол работает, но ему не хватает настройки. Добавим роль для LLM и входные данные от пользователя.

\begin{code}
ОПРЕДЕЛИ аналитик
<<
Ты опытный аналитик данных. Отвечаешь коротко и по делу.
>>
КОНЕЦ

ПОЛУЧИ запрос
<<
Вопрос пользователя для анализа.
>>
КОНЕЦ

ДУМАЙ ответ
  КАК $аналитик
  ЦЕЛЬ <<
  Ответь на вопрос пользователя.
  >>
  ВХОД <<
  $запрос
  >>
  РЕЗУЛЬТАТ
  * ответ: ТЕКСТ - ответ на вопрос
КОНЕЦ

ЖДИ
  НА ?ответ
КОНЕЦ

НАПИШИ
  КОМУ @пользователь
  <<
  $ответ.ответ
  >>
КОНЕЦ

ВЫХОД
\end{code}

Новые элементы:

\begin{itemize}[leftmargin=2em]
\item \ic{ОПРЕДЕЛИ аналитик}~--- создаёт именованное значение. Это мгновенный оператор: он не вызывает LLM, а просто сохраняет текст под именем \ic{\$аналитик}.

\item \ic{ПОЛУЧИ запрос}~--- блокирующий оператор. Протокол останавливается, пока среда не предоставит значение. Это может быть сообщение пользователя, параметр из конфигурации~--- протоколу всё равно, откуда.

\item \ic{КАК \$аналитик}~--- модификатор-оснащение в \ic{ДУМАЙ}. Говорит LLM, в какой роли действовать.

\item \ic{ВХОД <{<} ... >{>}}~--- модификатор-постановка. Условие задачи.
\end{itemize}

\subsubsection{Порядок модификаторов в ДУМАЙ}

Модификаторы \ic{ДУМАЙ} делятся на две группы, и группы не смешиваются:

\begin{enumerate}[leftmargin=2em]
\item \textbf{Оснащение} (чем и как):  \ic{ЧЕРЕЗ}, \ic{КАК}, \ic{ИСПОЛЬЗУЯ}
\item \textbf{Постановка} (что решать): \ic{ЦЕЛЬ}, \ic{ВХОД}, \ic{КОНТЕКСТ}, \ic{РЕЗУЛЬТАТ}
\end{enumerate}

Оснащение \textbf{всегда перед} постановкой. \ic{РЕЗУЛЬТАТ} \textbf{всегда последний}.

\subsection{Все сигилы}

В предыдущих примерах вы видели \ic{\$} и \ic{?}. Полный набор:

\begin{center}
\begin{tabular}{@{} c l l @{}}
\toprule
Сигил & Значение & Пример \\
\midrule
\ic{\$}  & Значение          & \ic{\$запрос}, \ic{\$план.summary} \\
\ic{?}  & Обещание          & \ic{?план}, \ic{?ответ} \\
\ic{@}  & Участник          & \ic{@тех\_аналитик}, \ic{@пользователь} \\
\ic{!}  & Инструмент        & \ic{!загрузить\_статью} \\
\ic{\#} & Адрес (канал)     & \ic{\#быстрые\_вопросы} \\
\ic{\~{}} & Поток            & \ic{\~{}решение} \\
\bottomrule
\end{tabular}
\end{center}

Участники и инструменты требуют предварительного объявления через \ic{УЧАСТНИКИ} и \ic{ИНСТРУМЕНТЫ}.

\subsection{Чек-лист: первый протокол}

Минимум для работающего протокола:

\begin{itemize}[leftmargin=2em]
\item Хотя бы один когнитивный шаг (\ic{ДУМАЙ}) или действие (\ic{ВЫПОЛНИ}, \ic{НАПИШИ}).
\item \ic{ЖДИ} перед использованием результата.
\item Явное завершение (\ic{ВЫХОД}).
\end{itemize}

Рекомендуется:

\begin{itemize}[leftmargin=2em]
\item \ic{УЧАСТНИКИ}~--- кто участвует в протоколе.
\item \ic{ИНСТРУМЕНТЫ}~--- какие инструменты доступны.
\item \ic{ОПРЕДЕЛИ}~--- роли и константы.
\end{itemize}


% ============================================================
% SECTION 2: ТРИ ЧАСТЫХ СЦЕНАРИЯ
% ============================================================

\section{Три частых сценария}

Большинство агентных задач укладываются в три паттерна. Освойте их~--- и вы закроете 80\% случаев.

\subsection{Классификация и маршрутизация}

Задача: входящий запрос нужно классифицировать и направить подходящему специалисту.

\begin{code}
ОПРЕДЕЛИ general_role
<<
Ты специалист по обслуживанию клиентов.
>>
КОНЕЦ

ОПРЕДЕЛИ technical_role
<<
Ты специалист технической поддержки.
Решаешь проблемы пошагово.
>>
КОНЕЦ

' --- Классификация ---

ДУМАЙ classification
  ЦЕЛЬ <<
  Классифицируй запрос клиента.
  >>
  ВХОД <<
  $message
  >>
  РЕЗУЛЬТАТ
  * type: ВЫБОР(general, technical) - тип запроса
КОНЕЦ

ЖДИ
  НА ?classification
КОНЕЦ

' --- Выбор роли ---

ОПРЕДЕЛИ current_role
$general_role
КОНЕЦ

ЕСЛИ $classification.type == "technical"
  УСТАНОВИ $current_role
  $technical_role
  КОНЕЦ
КОНЕЦ

' --- Ответ от выбранного специалиста ---

ДУМАЙ response
  КАК $current_role
  ВХОД <<
  $message
  >>
  РЕЗУЛЬТАТ
  * answer: ТЕКСТ - ответ клиенту
КОНЕЦ

ЖДИ
  НА ?response
КОНЕЦ

НАПИШИ
  КОМУ @пользователь
  <<
  $response.answer
  >>
КОНЕЦ

ВЫХОД
\end{code}

\subsubsection{Что здесь важно}

\begin{itemize}[leftmargin=2em]
\item \ic{ВЫБОР(general, technical)} в \ic{РЕЗУЛЬТАТ}~--- ограничивает LLM фиксированным набором значений. Именно это делает ветвление надёжным: \ic{ЕСЛИ} получает гарантированно корректное значение.

\item \ic{ЕСЛИ}~--- вычисляется средой исполнения \textbf{детерминированно}, без вызова LLM. Это управляющая логика, а не когнитивная.

\item Паттерн: \textbf{сначала думаем, чтобы классифицировать}~--- затем ветвимся детерминированно по результату.
\end{itemize}

\subsection{Параллельная экспертиза}

Задача: два эксперта параллельно анализируют код, техлид агрегирует результаты.

\begin{code}
ОПРЕДЕЛИ security_role
<<
Ты эксперт по безопасности кода.
>>
КОНЕЦ

ОПРЕДЕЛИ performance_role
<<
Ты эксперт по производительности кода.
>>
КОНЕЦ

ОПРЕДЕЛИ tech_lead_role
<<
Ты технический руководитель.
Подводишь итоги нескольких проверок.
>>
КОНЕЦ

' --- Два ревью запускаются подряд ---
' --- Оба работают параллельно ---

ДУМАЙ security_review
  КАК $security_role
  ВХОД <<
  $message
  >>
  РЕЗУЛЬТАТ
  * issues: ТЕКСТ - найденные уязвимости
  * risk: ВЫБОР(low, medium, high) - уровень риска
КОНЕЦ

ДУМАЙ performance_review
  КАК $performance_role
  ВХОД <<
  $message
  >>
  РЕЗУЛЬТАТ
  * issues: ТЕКСТ - найденные проблемы
  * impact: ВЫБОР(low, medium, high) - влияние
КОНЕЦ

' --- Ждём оба результата ---

ЖДИ
  НА ?security_review, ?performance_review
  РЕЖИМ ВСЕ
КОНЕЦ

' --- Агрегация ---

ДУМАЙ summary
  КАК $tech_lead_role
  ЦЕЛЬ <<
  Объедини результаты двух ревью в итог.
  >>
  КОНТЕКСТ <<
  Безопасность: $security_review
  Производительность: $performance_review
  >>
  РЕЗУЛЬТАТ
  * summary: ТЕКСТ - итоговое резюме
  * actions: ТЕКСТ - ключевые действия
КОНЕЦ

ЖДИ
  НА ?summary
КОНЕЦ

НАПИШИ
  КОМУ @пользователь
  <<
  $summary.summary

  Действия: $summary.actions
  >>
КОНЕЦ

ВЫХОД
\end{code}

\subsubsection{Что здесь важно}

\begin{itemize}[leftmargin=2em]
\item \textbf{Запускающие операторы не блокируют.} \ic{ДУМАЙ security\_review} создаёт обещание и тут же переходит к следующему \ic{ДУМАЙ performance\_review}. Оба работают параллельно.

\item \ic{ЖДИ НА ?security\_review, ?performance\_review РЕЖИМ ВСЕ}~--- явная точка синхронизации. Протокол ждёт, пока \textbf{оба} результата не будут готовы.

\item \ic{РЕЖИМ ЛЮБОЙ}~--- альтернатива: продолжить, как только готов хотя бы один результат.

\item Паттерн: \textbf{запустить несколько}, потом \textbf{подождать все}, потом \textbf{агрегировать}.
\end{itemize}

\subsection{Итеративное улучшение}

Задача: перевод текста с циклом проверки качества. Оценщик проверяет перевод, и если качество недостаточно~--- переводчик улучшает.

\begin{code}
ОПРЕДЕЛИ translator_role
<<
Ты эксперт по литературному переводу.
>>
КОНЕЦ

ОПРЕДЕЛИ evaluator_role
<<
Ты эксперт по оценке литературных переводов.
>>
КОНЕЦ

' --- Первый перевод ---

ДУМАЙ initial_translation
  КАК $translator_role
  ЦЕЛЬ <<
  Переведи текст, сохраняя тон и нюансы.
  >>
  ВХОД <<
  $message
  >>
  РЕЗУЛЬТАТ
  * text: ТЕКСТ - перевод
КОНЕЦ

ЖДИ
  НА ?initial_translation
КОНЕЦ

ОПРЕДЕЛИ current
$initial_translation.text
КОНЕЦ

ОПРЕДЕЛИ done
0
КОНЕЦ

' --- Цикл оценки и улучшения ---

ПОВТОРЯЙ ДО $done НЕ БОЛЕЕ 3

  ДУМАЙ evaluation
    КАК $evaluator_role
    ЦЕЛЬ <<
    Оцени качество перевода.
    >>
    ВХОД <<
    Оригинал: $message
    Перевод: $current
    >>
    РЕЗУЛЬТАТ
    * score: ЧИСЛО - качество от 1 до 10
    * issues: ТЕКСТ - проблемы
  КОНЕЦ

  ЖДИ
    НА ?evaluation
  КОНЕЦ

  ЕСЛИ $evaluation.score >= 8
    УСТАНОВИ $done
    1
    КОНЕЦ
  КОНЕЦ

  ЕСЛИ $done == 0
    ДУМАЙ improved
      КАК $translator_role
      ЦЕЛЬ <<
      Улучши перевод на основе замечаний.
      >>
      ВХОД <<
      Оригинал: $message
      Перевод: $current
      >>
      КОНТЕКСТ <<
      Проблемы: $evaluation.issues
      >>
      РЕЗУЛЬТАТ
      * text: ТЕКСТ - улучшенный перевод
    КОНЕЦ

    ЖДИ
      НА ?improved
    КОНЕЦ

    УСТАНОВИ $current
    $improved.text
    КОНЕЦ
  КОНЕЦ

КОНЕЦ

НАПИШИ
  КОМУ @пользователь
  <<
  $current
  >>
КОНЕЦ

ВЫХОД
\end{code}

\subsubsection{Что здесь важно}

\begin{itemize}[leftmargin=2em]
\item \ic{ПОВТОРЯЙ ДО \$done НЕ БОЛЕЕ 3}~--- цикл с обязательным лимитом. Это не стилистическое требование, а часть конструкции: каждая итерация с \ic{ДУМАЙ} расходует бюджет.

\item \ic{УСТАНОВИ} меняет существующее значение (а не \ic{ОПРЕДЕЛИ}, который создаёт новое).

\item Паттерн: \textbf{работник создаёт}, \textbf{оценщик проверяет}, \textbf{цикл повторяется до нужного качества}.
\end{itemize}

% ============================================================
% SECTION 3: ДУМАЙ ИЛИ НАПИШИ С ЖДАТЬ?
% ============================================================

\section{ДУМАЙ или НАПИШИ с ЖДАТЬ?}

Это центральный выбор при проектировании протокола. Оба способа дают результат~--- но природа результата разная.

\subsection{Два способа получить ответ}

\begin{center}
\begin{tabular}{@{} l p{5.5cm} p{5.5cm} @{}}
\toprule
& \textbf{ДУМАЙ} & \textbf{НАПИШИ с ЖДАТЬ} \\
\midrule
Кто работает & Текущий агент (LLM) & Другой участник \\
Результат определяется & Структурой \ic{РЕЗУЛЬТАТ} & Ответом участника \\
Контроль формата & Полный (типы, поля) & Зависит от участника \\
Когда использовать & Нужна когнитивная работа этого агента & Нужен ответ от другого агента или человека \\
\bottomrule
\end{tabular}
\end{center}

\subsection{Правило выбора}

\begin{itemize}[leftmargin=2em]
\item \textbf{Агент справится сам?} $\rightarrow$ \ic{ДУМАЙ}
\item \textbf{Нужен ответ от другого?} $\rightarrow$ \ic{НАПИШИ} с \ic{ЖДАТЬ}
\end{itemize}

Примеры:

\begin{center}
\begin{tabular}{@{} p{7cm} l @{}}
\toprule
Задача & Оператор \\
\midrule
Классифицировать запрос & \ic{ДУМАЙ} \\
Спросить у техника диагноз & \ic{НАПИШИ ЖДАТЬ ЛЮБОЙ} \\
Собрать выжимку из данных & \ic{ДУМАЙ} \\
Уточнить у пользователя детали & \ic{НАПИШИ ЖДАТЬ ЛЮБОЙ} \\
Попросить троих экспертов & \ic{НАПИШИ ЖДАТЬ ВСЕ} \\
\bottomrule
\end{tabular}
\end{center}

\subsection{НАПИШИ без ЖДАТЬ}

Если \ic{ЖДАТЬ} не указан и у \ic{НАПИШИ} нет имени результата~--- это отправка без ожидания: сообщение ушло, протокол продолжается.

Используйте для: финального ответа, уведомлений, статусных сообщений.

\begin{code}
' Финальный ответ --- не ждём
НАПИШИ
  КОМУ @пользователь
  <<
  Всё готово: $result.summary
  >>
КОНЕЦ
\end{code}

\subsection{Три политики ожидания}

\begin{center}
\begin{tabular}{@{} l p{8cm} @{}}
\toprule
Политика & Когда обещание разрешается \\
\midrule
\ic{ЖДАТЬ НЕТ}    & Сообщение отправлено (аналог отсутствия \ic{ЖДАТЬ}) \\
\ic{ЖДАТЬ ЛЮБОЙ}  & Получен хотя бы один ответ \\
\ic{ЖДАТЬ ВСЕ}    & Получены ответы от всех адресатов \\
\bottomrule
\end{tabular}
\end{center}

Пример запроса с ожиданием:

\begin{code}
НАПИШИ вопрос
  КОМУ @тех_аналитик
  ЖДАТЬ ЛЮБОЙ
  НЕ БОЛЕЕ 10м
  <<
  Проверь возможную причину сбоя.
  >>
КОНЕЦ

ЖДИ
  НА ?вопрос
КОНЕЦ

' Теперь $вопрос содержит ответ аналитика
\end{code}


% ============================================================
% SECTION 4: КАК НЕ ДЕЛАТЬ
% ============================================================

\section{Как не делать}

\subsection{Всё в одном ДУМАЙ}

Соблазн: описать весь рабочий процесс в одном гигантском промпте.

\begin{badcode}
' Плохо: вся логика спрятана в промпте
ДУМАЙ всё_сразу
  ЦЕЛЬ <<
  Классифицируй запрос, выбери специалиста,
  составь ответ, проверь качество, и если
  плохо --- переделай.
  >>
  РЕЗУЛЬТАТ
  * answer: ТЕКСТ
КОНЕЦ
\end{badcode}

Почему плохо:
\begin{itemize}[leftmargin=2em]
\item Логика ветвлений спрятана внутри промпта~--- теряется детерминизм.
\item Невозможно проверить промежуточные результаты.
\item Один шаг тратит весь бюджет.
\item При ошибке на любом этапе~--- всё заново.
\end{itemize}

Правильно: разбить на отдельные шаги, соединённые управляющей логикой (как в сценарии «Классификация и маршрутизация»).

\subsection{ДУМАЙ для детерминированного решения}

\begin{badcode}
' Плохо: LLM решает, больше ли score 80
ДУМАЙ check
  ЦЕЛЬ <<
  Определи, превышает ли оценка 80.
  >>
  ВХОД <<
  Оценка: $evaluation.score
  >>
  РЕЗУЛЬТАТ
  * exceeds: ФЛАГ
КОНЕЦ
\end{badcode}

Зачем тратить токены на то, что среда вычислит мгновенно?

\begin{code}
' Правильно: детерминированная проверка
ЕСЛИ $evaluation.score >= 80
  ...
КОНЕЦ
\end{code}

\ic{ЕСЛИ} вычисляется без LLM. Оставляйте \ic{ДУМАЙ} для задач, где нужно \textbf{понимание}, а не арифметика.

\subsection{Забыть ЖДИ перед использованием результата}

\begin{badcode}
' Плохо: $plan ещё не готов
ДУМАЙ plan
  ЦЕЛЬ <<
  Составь план.
  >>
  РЕЗУЛЬТАТ
  * steps: ТЕКСТ
КОНЕЦ

НАПИШИ
  КОМУ @пользователь
  <<
  План: $plan.steps
  >>
КОНЕЦ
\end{badcode}

После \ic{ДУМАЙ} существует \textbf{обещание} \ic{?plan}, но \textbf{значение} \ic{\$plan} ещё недоступно. Нужна явная синхронизация:

\begin{code}
ДУМАЙ plan
  ...
КОНЕЦ

ЖДИ
  НА ?plan
КОНЕЦ

НАПИШИ
  КОМУ @пользователь
  <<
  План: $plan.steps
  >>
КОНЕЦ
\end{code}

\subsection{Повторное ОПРЕДЕЛИ}

\begin{badcode}
' Плохо: ошибка подготовки
ОПРЕДЕЛИ counter
0
КОНЕЦ

ОПРЕДЕЛИ counter
1
КОНЕЦ
\end{badcode}

\ic{ОПРЕДЕЛИ} создаёт \textbf{новое} имя. Повторное \ic{ОПРЕДЕЛИ} с тем же именем~--- ошибка подготовки. Для изменения~--- \ic{УСТАНОВИ}:

\begin{code}
ОПРЕДЕЛИ counter
0
КОНЕЦ

УСТАНОВИ $counter
1
КОНЕЦ
\end{code}

\subsection{Цикл без лимита}

В COIL \ic{ПОВТОРЯЙ} без лимита невозможен~--- это часть конструкции. Зачем?

Каждая итерация с \ic{ДУМАЙ} внутри тратит токены. Без верхней границы протокол может исчерпать весь бюджет. Лимит~--- это не перестраховка, а осознанное проектное решение.


% ============================================================
% SECTION 5: ТИПИЧНЫЕ ОШИБКИ
% ============================================================

\section{Типичные ошибки}

\subsection{Порядок модификаторов в ДУМАЙ}

\begin{badcode}
' Ошибка подготовки: ЦЕЛЬ (постановка) перед КАК (оснащение)
ДУМАЙ plan
  ЦЕЛЬ <<
  Найди проблему.
  >>
  КАК $analyst_role
КОНЕЦ
\end{badcode}

Правило: \textbf{оснащение} (\ic{ЧЕРЕЗ}, \ic{КАК}, \ic{ИСПОЛЬЗУЯ}) \textbf{перед постановкой} (\ic{ЦЕЛЬ}, \ic{ВХОД}, \ic{КОНТЕКСТ}, \ic{РЕЗУЛЬТАТ}).

\begin{code}
ДУМАЙ plan
  КАК $analyst_role
  ЦЕЛЬ <<
  Найди проблему.
  >>
КОНЕЦ
\end{code}

\subsection{РЕЗУЛЬТАТ не последний}

\begin{badcode}
' Ошибка: КОНТЕКСТ после РЕЗУЛЬТАТ
ДУМАЙ analysis
  РЕЗУЛЬТАТ
  * summary: ТЕКСТ
  КОНТЕКСТ <<
  Дополнительные данные.
  >>
КОНЕЦ
\end{badcode}

\ic{РЕЗУЛЬТАТ} \textbf{обязан} быть последним модификатором в \ic{ДУМАЙ}.

\subsection{Шаблон <{<} >{>} в ВЫПОЛНИ}

\begin{badcode}
' Ошибка подготовки: шаблон запрещён в ВЫПОЛНИ
ВЫПОЛНИ article
  ИСПОЛЬЗУЯ !загрузить_статью
  <<
  Загрузи статью по адресу $url
  >>
КОНЕЦ
\end{badcode}

\ic{ВЫПОЛНИ}~--- вызов инструмента. Инструменты принимают структурированные аргументы, а не свободный текст:

\begin{code}
ВЫПОЛНИ article
  ИСПОЛЬЗУЯ !загрузить_статью
  - url: $url
  - format: "markdown"
КОНЕЦ
\end{code}

Шаблоны \ic{<{<} >{>}} предназначены для естественного языка (\ic{ДУМАЙ}, \ic{НАПИШИ}). Аргументы \ic{- key: value}~--- для вызовов инструментов (\ic{ВЫПОЛНИ}).

\subsection{Необъявленные участники и инструменты}

\begin{badcode}
' Ошибка подготовки: @аналитик не объявлен
НАПИШИ
  КОМУ @аналитик
  <<
  Проверь данные.
  >>
КОНЕЦ
\end{badcode}

Все участники должны быть объявлены через \ic{УЧАСТНИКИ}, все инструменты~--- через \ic{ИНСТРУМЕНТЫ}:

\begin{code}
УЧАСТНИКИ аналитик

НАПИШИ
  КОМУ @аналитик
  <<
  Проверь данные.
  >>
КОНЕЦ
\end{code}

\subsection{Путаница \$name и ?name}

\begin{badcode}
' Ошибка: ЖДИ ожидает обещание (?), а не значение ($)
ЖДИ
  НА $plan
КОНЕЦ
\end{badcode}

\ic{ЖДИ НА} принимает \textbf{обещания} (\ic{?name}). \ic{\$name}~--- это уже разрешённое значение, ждать его бессмысленно.

\begin{code}
ЖДИ
  НА ?plan
КОНЕЦ
\end{code}


% ============================================================
% SECTION 6: ЧТО ДАЛЬШЕ
% ============================================================

\section{Что дальше}

\subsection{Extended-операторы}

В примерах выше вы уже использовали \ic{ЕСЛИ} и \ic{ПОВТОРЯЙ}. Они относятся к Extended-набору, который также включает:

\begin{itemize}[leftmargin=2em]
\item \ic{КАЖДЫЙ}~--- итерация по элементам списка.
\item \ic{СОБЕРИ}~--- агрегация набора результатов.
\item \ic{СИГНАЛ}~--- отправка данных в существующий поток.
\end{itemize}

Extended-операторы стабильны по идее, но их точная форма может уточняться.

\subsection{Инструменты и ВЫПОЛНИ}

\ic{ВЫПОЛНИ} вызывает внешние инструменты: загрузку данных, API, интерактивные сессии. Инструмент создаёт и обещание \ic{?name}, и поток \ic{\~{}name} (для интерактивных сессий).

\begin{code}
ИНСТРУМЕНТЫ загрузить_статью

ВЫПОЛНИ article
  ИСПОЛЬЗУЯ !загрузить_статью
  - url: $url
КОНЕЦ
\end{code}

\subsection{Потоки и события}

Каждый запускающий оператор создаёт не только обещание (\ic{?name}), но и поток (\ic{\~{}name}). Поток~--- это двусторонний канал для отправки и получения данных. Через \ic{СИГНАЛ} можно дослать данные в работающий шаг:

\begin{code}
СИГНАЛ ~решение
  <<
  Клиент добавил новую информацию.
  >>
КОНЕЦ
\end{code}

\subsection{Где искать примеры}

\begin{itemize}[leftmargin=2em]
\item \ic{examples/hello-world.md}~--- минимальный протокол
\item \ic{examples/patterns/routing.coil}~--- классификация и маршрутизация
\item \ic{examples/patterns/parallelization.coil}~--- параллельная экспертиза
\item \ic{examples/patterns/evaluator-optimizer.coil}~--- итеративное улучшение
\item \ic{examples/patterns/orchestrator-workers.coil}~--- планирование и параллельные воркеры
\item \ic{examples/patterns/prompt-chaining.coil}~--- цепочка промптов
\end{itemize}

Полная спецификация~--- в директории \ic{spec/}.

\end{document}
